\section{Introduction}
\label{sec:intro}

% The Vehicle Routing Problem (VRP) is a classic optimization problem that has been revisited due to its significant real-world applications and the new opportunities brought by emerging technologies such as artificial intelligence and the Internet of Things (IoT). Efficiently managing a fleet of vehicles in real time to service multiple locations under various constraints—such as delivery time windows, vehicle capacities, and route costs and risks — has the potential to greatly enhance logistics, reduce fuel consumption, improve customer satisfaction, and safety. However, solving the VRP is complex and is often classified as NP-hard, indicating that the time required to find an optimal solution increases significantly as the problem size grows, making it computationally challenging for large instances,  especially under short time constraints.

% To address this complexity, many large-scale technology companies have developed powerful cloud-based APIs aimed at providing route optimization  as a service. These APIs often combine sophisticated heuristics algorithms with scalable cloud resources to tackle the VRP across diverse industries, from e-commerce to transportation. Solutions from companies like Google\footnote{\url{https://developers.google.com/maps/documentation/route-optimization}}, Microsoft\footnote{\url{https://learn.microsoft.com/en-us/rest/api/maps/route}}, and others offer varying degrees of optimization based on specific parameters such as cost-efficiency, speed, or the ability to handle large datasets. However, given the diversity of API offerings, companies often face a difficult decision: which service best meets their needs in terms of functionality, performance, and cost-effectiveness?

% In this context, selecting the most appropriate API is not a trivial task. The need for flexibility and customization across varying use cases—whether for optimizing cost, increasing input size, or adapting to specific business constraints—complicates the decision further. Enterprises may find themselves limited by the features or pricing structures of a single API, or they may need to constantly switch between services to best suit different routing scenarios.

% This paper introduces RouteBastion, a Software as a Service (SaaS) middleware solution designed to unifying multiple Vehicle Routing Problem APIs into a single platform. RouteBastion abstracts the complexity of comparing and integrating different APIs by providing a seamless, centralized service where users can dynamically select the most suitable API for their needs. The platform’s architecture is built on principles of modularity and scalability, ensuring that it can accommodate the constantly evolving landscape of route optimization services. This SaaS platform will offer more than just API aggregation; it will incorporate intelligent decision-making mechanisms, allowing users to optimize route planning not only based on individual API strengths but also on real-time operational requirements, such as route complexity, input size, or pricing considerations.

% In the remainder of this paper, we first present the theoretical framework that underpins our research, providing the necessary context and concepts for understanding the Vehicle Routing Problem and its implications (Section~\ref{sec:problem_analysis}). 
% %Following this, we conduct a thorough problem analysis to identify the key challenges and requirements associated with existing Vehicle Routing Problem solutions. 
% Next, we introduce our architecture proposal, detailing the design and components of RouteBastion ((Section~\ref{sec:routebastion}). This is followed by a evaluation of the proposed architecture to assess its effectiveness and potential impact (Section~\ref{sec:evaluation}). Finally, we conclude with a summary of our findings and implications for future research~(Section~\ref{sec:conclusion}).

% To address this complexity, many large-scale technology companies have developed powerful cloud-based APIs aimed at providing route optimization software as a service (SaaS). These APIs often combine sophisticated heuristics algorithms with scalable cloud resources to tackle the VRP across diverse industries, such as e-commerce and transportation. However, selecting the most appropriate API is not a trivial task, as companies often face the need for flexibility and customization across varying use cases \cite{Brokering2019JPDC,cloudbrokersms2024ieeetransactions,MCDM2022ieeeaccess}.

The Vehicle Routing Problem (VRP) is a complex optimization problem that has gained significant attention due to its real-world applications and the opportunities brought by emerging technologies like artificial intelligence and IoT\cite{10379532,SHAHIN2024104437}. Efficiently managing a fleet of vehicles in real-time under various constraints can enhance logistics and customer satisfaction. However, solving the VRP is classified as NP-hard, making it computationally challenging for large instances, especially under short time constraints\cite{10379532}.

To address the complexity of route optimization, large-scale technology companies provide cloud-based APIs solutions, using advanced heuristic algorithms and scalable cloud infrastructure. The selection of an appropriate API creates opportunities for specialized brokers to manage the arbitration process, orchestrating services to ensure they meet specific business requirements \cite{Brokering2019JPDC,cloudbrokersms2024ieeetransactions}. However, optimizing service-level agreements (SLAs) while balancing performance, cost, and customization remains a challenging task \cite{MCDM2022ieeeaccess}.

This paper introduces RouteBastion, a Software-as-a-service (SaaS) middleware solution designed to unify multiple VRP APIs into a single platform. RouteBastion abstracts the complexity of comparing and integrating different APIs by providing a seamless, centralized service where users can dynamically select the most suitable API for their needs. The platform's architecture is built on principles of modularity and scalability, ensuring it can accommodate the constantly evolving landscape of route optimization services.

The paper presents a problem analysis (Section~\ref{sec:problem_analysis}), introduces the architecture proposal (Section~\ref{sec:routebastion}), evaluates the proposed architecture (Section~\ref{sec:evaluation}), and concludes with a summary of findings and implications for future research (Section~\ref{sec:conclusion}).
